% ============================================================
%  14_Results.tex – BULGULAR
% ============================================================
\section{BULGULAR}
\label{sec:bulgular}

Bu bölümde, geliştirilen Myora platformunun modül bazında işlevselliği, gerçekleştirilen kod metrikleri, yapay zekâ analiz çıktıları ve veritabanı istatistikleri sunulmaktadır.

\subsection{Gerçekleştirilen Modüller ve Durumları}

Proje kapsamında planlanan 13 temel modülün geliştirme durumları, kaynak koddan üretilen analiz sonuçlarına göre Tablo~\ref{tab:module_status}'de özetlenmektedir.

\begin{table}[H]
\centering
\caption{Proje modüllerinin geliştirme durumları (kod tabanından doğrulanmıştır)}
\label{tab:module_status}
\small
\begin{tabular}{|c|l|c|p{0.45\textwidth}|}
\hline
\textbf{\#} & \textbf{Modül} & \textbf{Durum} & \textbf{Kod Konumu} \\
\hline
1  & AI Sağlık Profili         & Tam     & \texttt{HealthProfile.tsx} + \texttt{generate-health-profile} edge fn \\
2  & RAG Chatbot               & Tam     & \texttt{ExpertChat.tsx} + \texttt{generate-chat-response} (8 aşama) \\
3  & Kan Tahlili Analizi       & Tam     & \texttt{BloodTestAnalysis.tsx} + \texttt{analyze-health-data} (\texttt{blood-test}) \\
4  & Yemek Fotoğraf Analizi    & Tam     & \texttt{FoodCapture.tsx} + \texttt{analyze-health-data} (\texttt{food-image}) \\
5  & Takviye Gıda Mağazası     & Tam     & \texttt{SupplementStore.tsx}, \texttt{SupplementTracker.tsx} \\
6  & Aile Takibi               & Tam     & \texttt{FamilyTracking.tsx} + \texttt{family\_*} tabloları \\
7  & Topluluk                  & Tam     & \texttt{Community.tsx} + \texttt{community\_*} tabloları \\
8  & Giyilebilir Cihaz Entegr. & Kısmi   & \texttt{WearableDevices.tsx}, \texttt{WearableDataDashboard.tsx} (katalog tabanlı) \\
9  & Tıbbi Fotoğraf Analizi    & Tam     & \texttt{MedicalPhotoAnalysis.tsx} + 5 fotoğraf tipi \\
10 & Sesli Günlük              & Kısmi   & \texttt{VoiceLogger.tsx} + mock transkripsiyon (Whisper henüz değil) \\
11 & Aralıklı Oruç             & Tam     & \texttt{IntervalFasting.tsx} + \texttt{fasting\_logs} \\
12 & Hedefler \& Başarımlar    & Tam     & \texttt{AchievementsBadges.tsx} + \texttt{user\_goals} / \texttt{achievements} \\
13 & Akıllı Bildirimler        & Tam     & \texttt{NotificationSettings.tsx} + Capacitor LocalNotifications \\
\hline
\end{tabular}
\\[4pt]
\footnotesize{Tam: Backend ve frontend tam entegre, Kısmi: Backend hazır / frontend kısmen entegre}
\end{table}

\subsection{Kod Metrikleri}

Tablo~\ref{tab:code_metrics} repository'den çıkarılan ham metrikleri özetlemektedir.

\begin{table}[H]
\centering
\caption{Kaynak kod metrikleri}
\label{tab:code_metrics}
\begin{tabular}{|l|c|}
\hline
\textbf{Metrik} & \textbf{Değer} \\
\hline
React üst-düzey bileşeni (\texttt{client/src/components})       & 41 \\
shadcn/ui temel bileşeni (\texttt{components/ui})              & 53 \\
Sayfa kabuğu (\texttt{client/src/pages})                       & 5 \\
React Context sağlayıcısı                                     & 3 \\
Özel React kancası (custom hook)                              & 7 \\
İstemci API modülü (\texttt{client/src/lib/api})              & 8 + ortak/types \\
Toplam istemci API satırı (TypeScript)                         & 2954 \\
Supabase Edge Function                                        & 3 (1340 satır) \\
RAG kütüphane modülü (\texttt{generate-chat-response/lib})     & 5 (843 satır) \\
SQL migrasyon dosyası                                         & 14 (1219 satır) \\
PostgreSQL tablosu                                            & 33 \\
Paylaşılan tip dosyası (\texttt{shared/schema.ts})            & 646 satır \\
RAG ingest betiği (\texttt{scripts/rag})                       & 2 (1124 satır) \\
\hline
\end{tabular}
\end{table}

\subsection{Yapay Zekâ Analiz Sonuçları}

\subsubsection{Yemek Fotoğraf Analizi}

GPT-4o vision modeli, kamera ile çekilen yemek fotoğraflarını analiz ederek yapılandırılmış JSON yanıt döndürmektedir. Sistem ayrıca AI'ın güvenilirlik oranını (\texttt{confidence}) raporlamakta, OpenAI çağrısı başarısız olursa kural-tabanlı yedek (fallback) bir çıktı sunmaktadır (Şekil~\ref{lst:food_json}).

\begin{figure}[H]
\centering
\begin{minipage}{0.85\textwidth}
\small\ttfamily
\noindent\{ \\
\hspace*{1em}"food": "Mercimek Çorbası", \\
\hspace*{1em}"calories": 180, \\
\hspace*{1em}"servingSize": "1 kase (250 mL)", \\
\hspace*{1em}"confidence": 92, \\
\hspace*{1em}"macros": \{ \\
\hspace*{2em}"protein": \{ "amount": 12, "percentage": 27 \}, \\
\hspace*{2em}"carbs":   \{ "amount": 28, "percentage": 62 \}, \\
\hspace*{2em}"fat":     \{ "amount":  3, "percentage": 15 \}, \\
\hspace*{2em}"fiber":   \{ "amount":  8, "percentage": 32 \} \\
\hspace*{1em}\}, \\
\hspace*{1em}"healthNotes": ["Yüksek lif içeriği", "Demir kaynağı"] \\
\}
\end{minipage}
\caption{\texttt{analyze-health-data} edge fonksiyonunun \texttt{food-image} alt-tipi için JSON çıktı örneği}
\label{lst:food_json}
\end{figure}

\subsubsection{Kan Tahlili Analizi}

\texttt{blood-test} alt-tipi, OCR ile dijitalleştirilen kan tahlili metinlerinden belirteçleri çıkarmakta ve her birini referans aralığına göre \texttt{normal} / \texttt{low} / \texttt{high} / \texttt{critical} olarak sınıflandırmaktadır. Sonuçlar \texttt{blood\_test\_results} tablosunda \texttt{(test\_id, marker\_name, value, unit, reference\_range, status, category, ai\_interpretation)} alanlarıyla saklanır.

\subsubsection{Tıbbi Fotoğraf Analizi}

\texttt{medical-photo} alt-tipi 5 farklı fotoğraf türü için ön değerlendirme üretir (Tablo~\ref{tab:medical_photo}). Her tür için \texttt{analysis} nesnesi renk, kıvam, gözlemler, endişeler, öneriler ve aciliyet alanlarını içerir.

\begin{table}[H]
\centering
\caption{Tıbbi fotoğraf analizi destek tipleri}
\label{tab:medical_photo}
\begin{tabular}{|l|l|}
\hline
\textbf{Tip} & \textbf{Tipik AI çıktısı} \\
\hline
\texttt{urine}  & Renk, berraklık, hidrasyon değerlendirmesi \\
\texttt{stool}  & Renk, kıvam, sindirim ipuçları \\
\texttt{tongue} & Renk, kaplama, dolaşım ipuçları \\
\texttt{eye}    & Skleral renk, kanlanma, anemi/sarılık ipucu \\
\texttt{skin}   & Renk değişimi, döküntü tipi, aciliyet \\
\hline
\end{tabular}
\end{table}

\subsubsection{Sağlık Profili Oluşturma}

\texttt{generate-health-profile} edge fonksiyonu kullanıcı verilerini tek bir kapsamlı profilde birleştirmektedir. Çıktıda BMI değeri, Mifflin--St Jeor formülüyle BMR (\eqref{eq:bmr_male}, \eqref{eq:bmr_female}), aktivite çarpanıyla TDEE ve 0--100 arası sağlık skoru bulunmaktadır.

\begin{equation}
\label{eq:bmr_male}
\mathrm{BMR}_{\text{erkek}} = 10 \cdot W + 6{,}25 \cdot H - 5 \cdot A + 5
\end{equation}

\begin{equation}
\label{eq:bmr_female}
\mathrm{BMR}_{\text{kadın}} = 10 \cdot W + 6{,}25 \cdot H - 5 \cdot A - 161
\end{equation}

\begin{equation}
\label{eq:tdee}
\mathrm{TDEE} = \mathrm{BMR} \cdot k_{\text{aktivite}}
\end{equation}

Burada $W$ kütle (kg), $H$ boy (cm), $A$ yaş (yıl) ve $k_{\text{aktivite}} \in \{1{,}2;\ 1{,}375;\ 1{,}55;\ 1{,}725;\ 1{,}9\}$ aktivite seviyesi çarpanıdır.

\subsection{Veritabanı İstatistikleri}

\begin{table}[H]
\centering
\caption{Veritabanı şema istatistikleri}
\label{tab:db_stats}
\begin{tabular}{|l|c|}
\hline
\textbf{Metrik} & \textbf{Değer} \\
\hline
Toplam tablo sayısı                       & 33 \\
UUID birincil anahtar kullanan tablo      & 33 \\
JSONB alanı içeren tablo                  & 9+ \\
Tetikleyici (trigger) sayısı              & 3+ (\texttt{handle\_new\_user\_profile}, \texttt{set\_*\_updated\_at}) \\
PostgreSQL fonksiyonu                     & 4+ (\texttt{match\_source\_chunks}, \texttt{handle\_new\_user\_profile}, \texttt{set\_profiles\_updated\_at}, \texttt{handle\_new\_user\_streak}) \\
RLS aktif olan tablo                      & 28+ \\
RLS politikası sayısı                     & 70+ \\
İndeks (B-tree, \texttt{ivfflat})         & 15+ \\
pgvector boyutu                           & 1536 \\
\hline
\end{tabular}
\end{table}

\subsection{API Modülleri ve Test Yaklaşımı}

İstemci API katmanı, 8 modüler dosyaya ayrılmış olup tüm sorgular Supabase JS SDK üzerinden gerçekleştirilmektedir. Tüm korumalı sorgular RLS politikaları sayesinde kimlik doğrulanmış kullanıcının kendi verisine erişimine sınırlandırılmıştır. Edge Function'lar Zod tabanlı doğrulama yerine \texttt{String()}/\texttt{Number()} dönüşümleri ile defansif programlama yaklaşımı kullanmaktadır; OpenAI çağrısının başarısız olduğu her durumda anlamlı bir fallback yanıt üretilmektedir.

\subsection{Performans ve Erişilebilirlik İyileştirmeleri}

\texttt{App.tsx} içinde tüm sayfa-düzeyinde bileşenler \texttt{React.lazy()} ile tembel yüklendiğinden ilk yükleme paket boyutu küçük tutulmuştur. \texttt{ErrorBoundary} sınıf bileşeni rota seviyesinde kullanıcı dostu bir hata UI sunar. Recharts grafikler ayrı sarmalayıcı bileşenlerle (\texttt{components/charts/}) tembel yüklenir. \texttt{LanguageContext} \texttt{toLocaleDateString}, \texttt{toLocaleTimeString} ve \texttt{toLocaleString} çağrılarına \texttt{locale} parametresini iletmektedir; böylece tarih ve sayı biçimlendirmesi kullanıcı dilinde doğru görüntülenir.

\subsection{Mobil Paketleme}

\texttt{capacitor.config.ts} dosyasında \texttt{appId=com.myora.app}, \texttt{webDir=dist/public} olarak tanımlanmıştır. \texttt{npm run build} \texttt{$\rightarrow$} \texttt{npx cap sync android} ardışık komutu ile Android Studio projesi güncellenmektedir. \texttt{MainActivity.java} içinde özel WebView kurulumu (kamera/mikrofon izinleri, dosya seçici) \texttt{onStart()} içinde gerçekleştirilmektedir.
